\chapter{Epiglottitis}
\index{Epiglottitis}
\label{throat:Epiglottitis}

\section{Overview}
Epiglottitis (supraglottitis) is a rapidly progressive, life-threatening cellulitis and swelling of the epiglottis and surrounding supraglottic structures (arytenoids, aryepiglottic folds), leading to catastrophic upper airway obstruction. Primary bacterial invasion causes massive supraglottic swelling, narrowing the laryngeal inlet.
\section{Causes}
\section{Risk Factors}

\begin{itemize}[noitemsep]
  \item Genetic predisposition and family history
  \item Advancing age
  \item Lifestyle factors (diet, exercise, smoking, alcohol)
  \item Environmental exposures
  \item Underlying medical conditions
  \item Medication use
\end{itemize}
\section{Signs and Symptoms}

\subsection{Common symptoms}
\begin{itemize}[noitemsep]
  \item \item Clinical features present with acute onset of high fever, severe sore throat, odynophagia, muffled 'hot potato' voice, drooling, and respiratory distress. Patients characteristically assume a 'tripod position' (sitting upright, leaning forward, neck hyperextended, chin thrust forward) to optimize airway caliber. Direct visualization with a tongue depressor is strictly contraindicated due to risk of precipitating complete laryngospasm. Lateral neck X-ray reveals a swollen epiglottis ('thumbprint sign'). Management mandates immediate airway stabilization in an operating room via endotracheal intubation or tracheostomy, followed by parenteral broad-spectrum antibiotics (ceftriaxone plus vancomycin). Pain or discomfort in the affected area    \item Pain or discomfort in the affected area
  \end{itemize}

\subsection{Emergency warning signs}
\begin{itemize}[noitemsep]
  \item Severe or worsening symptoms of epiglottitis require immediate medical evaluation
\end{itemize}
\section{Progression and Stages}
The condition may progress through different stages of severity over time. Early stages often involve mild or subtle symptoms that may be overlooked. Moderate stages involve more noticeable symptoms affecting daily function. Advanced stages may involve significant impairment and complications.
\section{Complications}
\begin{itemize}[noitemsep]
  \item Disease progression leading to worsening symptoms
  \item Secondary infections or injuries
  \item Side effects of treatment
  \item Reduced quality of life
  \item Emotional and psychological impact
\end{itemize}
\section{Diagnosis}
Comprehensive medical history and physical examination Laboratory tests to assess organ function and rule out other conditions Imaging studies to visualize affected structures Specialized testing depending on the affected system Consultation with relevant specialists Monitoring of disease progression over time

\begin{itemize}[noitemsep]
  \item Comprehensive medical history and physical examination
  \item Laboratory tests to assess organ function
  \item Imaging studies to visualize affected structures
  \item Specialized testing depending on the affected system
  \item Consultation with relevant specialists
\end{itemize}
\section{Prevention}
Historically caused by Haemophilus influenzae type b (Hib) in young children, widespread Hib vaccination has shifted epidemiology toward older adults and non-typeable H. influenzae, Streptococcus pneumoniae, Staphylococcus aureus, and Beta-hemolytic streptococci.

\begin{itemize}[noitemsep]
  \item Maintain a healthy lifestyle with balanced diet and exercise
  \item Avoid known risk factors
  \item Attend regular medical check-ups for early detection
  \item Follow recommended screening guidelines
\end{itemize}
\section{Treatment Options}
Medications to manage symptoms and address underlying mechanisms Lifestyle modifications and self-care measures Physical therapy and rehabilitation Surgical intervention when indicated Regular monitoring and follow-up care Patient education and support services

\begin{itemize}[noitemsep]
  \item Medications to manage symptoms and address underlying mechanisms
  \item Lifestyle modifications and self-care measures
  \item Physical therapy and rehabilitation
  \item Surgical intervention when indicated
  \item Regular monitoring and follow-up care
\end{itemize}
\section{Living with It}
\begin{itemize}[noitemsep]
  \item Follow your treatment plan as prescribed
  \item Maintain regular follow-up with your healthcare provider
  \item Make necessary lifestyle adjustments
  \item Seek support from family, friends, or support groups
  \item Stay informed about your condition
\end{itemize}
\section{Outlook}
The outlook varies depending on the specific condition, severity at diagnosis, and response to treatment. Early intervention generally leads to better outcomes. Many conditions can be effectively managed with appropriate treatment. Regular follow-up care is important for monitoring and treatment adjustment.
